\documentclass[../../main.tex]{subfiles}
\begin{document}
\begin{oframed}
  \subfile{problem}
\end{oframed}

{\bf [解答]}

\begin{figure}[htb]
  \centering
  \begin{tikzpicture}[>=stealth, every node/.style={font=\small}]
    % --- Panel 1: 四角錐と切断面 ---
    \begin{scope}[shift={(0,0)}]
      % 頂点座標
      \coordinate (V) at (0.3, 3.2);
      \coordinate (A) at (-1.8, -0.6);
      \coordinate (B) at (0.8, -0.9);
      \coordinate (C) at (2.4, 0.1);
      \coordinate (D) at (-0.2, 0.4);

      % 中点 M, N
      \coordinate (M) at ($ (A)!0.5!(D) $);
      \coordinate (N) at ($ (B)!0.5!(C) $);

      % 分点 P, Q, R
      \coordinate (P) at ($ (V)!0.75!(A) $);
      \coordinate (Q) at ($ (V)!0.75!(D) $);
      \coordinate (R) at ($ (P)!0.5!(Q) $);

      % 切断面 PBCQ のハッチング
      \fill[pattern=north east lines, pattern color=blue!45]
      (P) -- (B) -- (C) -- (Q) -- cycle;

      % 奥の隠れ線
      \draw[thick, dashed, gray!70!black] (A) -- (D) -- (C);
      \draw[thick, dashed, gray!70!black] (V) -- (D);
      \draw[dashed, gray!70!black] (M) -- (N);
      \draw[dashed, gray!70!black] (V) -- (M);

      % 手前の稜線
      \draw[line width=1.1pt, gray!85!black] (V) -- (A) -- (B) -- (C) -- (V);
      \draw[line width=1.1pt, gray!85!black] (V) -- (B);

      % 切断面の輪郭線
      \draw[line width=1.2pt, blue!85!black] (P) -- (B) -- (C);
      \draw[line width=1.2pt, blue!85!black, dashed] (C) -- (Q) -- (P);
      \draw[line width=1.1pt, red!85!black, dashed] (R) -- (N);

      % 頂点ノード
      \node[above=2pt] at (V) {$V$};
      \node[below left=1pt] at (A) {$A$};
      \node[below=2pt] at (B) {$B$};
      \node[right=2pt] at (C) {$C$};
      \node[above right=1pt, gray!70!black] at (D) {$D$};
      \node[left=2pt, blue!85!black] at (P) {$P$};
      \node[above right=1pt, blue!85!black] at (Q) {$Q$};
      \node[below left=1pt, gray!80!black] at (M) {$M$};
      \node[below right=1pt, gray!80!black] at (N) {$N$};
      \node[above left=1pt, red!85!black] at (R) {$R$};

      \node[below=1.2cm] at (0.3, 0) {(a) 四角錐と切断面 $PBCQ$};
    \end{scope}

    % --- Panel 2: 切断三角形 VMN ---
    \begin{scope}[shift={(5.2, 0)}]
      \coordinate (M2) at (0, -0.6);
      \coordinate (N2) at (3.0, -0.6);
      \coordinate (V2) at (1.5, 3.2);
      \coordinate (R2) at ($ (V2)!0.75!(M2) $);

      \draw[line width=1.1pt, gray!85!black] (M2) -- (V2) -- (N2) -- cycle;
      \draw[line width=1.3pt, red!85!black] (R2) -- (N2);

      % 点
      \node[circle, fill=gray!80!black, inner sep=1.2pt] at (M2) {};
      \node[below left=1pt] at (M2) {$M$};
      \node[circle, fill=gray!80!black, inner sep=1.2pt] at (N2) {};
      \node[below right=1pt] at (N2) {$N$};
      \node[circle, fill=gray!80!black, inner sep=1.2pt] at (V2) {};
      \node[above=2pt] at (V2) {$V$};
      \node[circle, fill=red!85!black, inner sep=1.2pt] at (R2) {};
      \node[above left=1pt, red!85!black] at (R2) {$R$};

      % 寸法表示
      \node[below=2pt] at ($ (M2)!0.5!(N2) $) {$20$};
      \node[above right=1pt] at ($ (V2)!0.5!(N2) $) {$40$};
      \node[above left=1pt] at ($ (V2)!0.5!(R2) $) {$30$};
      \node[below left=1pt] at ($ (R2)!0.5!(M2) $) {$10$};
      \node[above=2pt, red!85!black] at ($ (R2)!0.5!(N2) $) {$20$};

      \node[below=1.2cm] at (1.5, 0) {(b) 切断面 $\triangle VMN$};
    \end{scope}
  \end{tikzpicture}
  \caption{四角錐の切断と $\triangle VMN$ の断面}
  \label{fig:pyramid_cross_section}
\end{figure}

辺$VD$上に$VQ:QD=3:1$なる点$Q$をとると，切り口は台形$PBCQ$となる．相似から
$PQ=\dfrac{3}{4}AB=15 \,\mathrm{cm}$である．また$AD$，$BC$の中点を$M$，$N$として
$\triangle VMN$で立体を切断すると右のようになる．ただし$R$は$PQ$の中点である．
ここで$\triangle VMN$に余弦定理を用いて
\begin{align}
  \cos \angle RMN=\dfrac{1}{4}
\end{align}
だから，$\triangle MNR$に余弦定理を用いて
\begin{align}
  RN=\sqrt{(10)^2+(20)^2-100}=20 \,\mathrm{cm}
\end{align}
である．対称性から$RN$は台形の高さであるから，求める面積は
\begin{align}
  \dfrac{1}{2}|RN|(|PQ|+|BC|)&=\dfrac{1}{2}20(15+20) \\
  &=350 \,\mathrm{cm^2}
\end{align}
である．$\cdots$(答)

\newpage
\end{document}
